% GENERATED FILE. Edit canonical lessons, then run npm run generate:books.
% canonical-source-hash: fnv1a64:ef3351b687812b63
% canonical-lessons: PA-C325-sudharna, PA-C325-vigarna, PA-C325-parhez, PA-C325-chabauna, PA-C325-pir

\chapter{Getting Better}
\label{ch:pa-a2-getting-better}
\hlchaptermodality{325}

\begin{chapteropening}
\textbf{By the end of this chapter:} \emph{I can say whether someone is getting better or worse.}

Complete the last lesson of chapter 325: I can say whether someone is getting better or worse.
\end{chapteropening}

\section[\texorpdfstring{\pa{ਸੁਧਰਨਾ} (sudharnā)}{ਸੁਧਰਨਾ (sudharnā)}]{\pa{ਸੁਧਰਨਾ} (sudharnā) — to improve}

\label{lesson:PA-C325-sudharna}



\begin{quote}
Before the new one: say the Punjabi for sweat, then the Punjabi for a dizzy spell.
\end{quote}

\subsection*{You'll want to know: \pa{ਸੁਧਰਨਾ}}
\textbf{\pa{ਸੁਧਰਨਾ}} — \emph{sudharnā} — \textquotedblleft{}to improve\textquotedblright{}.

\begin{quote}
\pa{ਦਵਾਈ} \pa{ਤੋਂ} \pa{ਬਾਅਦ} \pa{ਦਾਦੀ} \pa{ਦੀ} \pa{ਸਿਹਤ} \pa{ਸੁਧਰੀ।}
\end{quote}

\begin{quote}
\emph{davāī ton bāad dādī dī sihat sudhrī.}
\end{quote}

\begin{quote}
After the medicine my grandmother's health got better.
\end{quote}

Past (he …): \pa{ਉਹ} \pa{ਇੱਕ} \pa{ਹਫ਼ਤੇ} \pa{ਵਿੱਚ} \pa{ਸੁਧਰਿਆ}, \emph{oh ikk hafte vich sudhriā}.

\begin{cousinweb}
\textbf{\pa{ਸੁਧਰਨਾ}} (\emph{sudharnā}) is inherited from a reconstructed \emph{*śuddhakāra}, \textquotedblleft{}making clean\textquotedblright{}; CDIAL lists Punjabi \emph{sudharnā}, \textquotedblleft{}to repair\textquotedblright{}, under it (CDIAL 12521).
\end{cousinweb}

\begin{grammarlens}[title={What you've built}]
One more word for this chapter.
\end{grammarlens}

\subsection*{Your turn}
\begin{itemize}
\raggedright
  \item \emph{Say it:} \emph{sudharnā}
  \item \emph{Say it:} \emph{sudharnā}, once more
  \item \emph{Say it:} \emph{davāī ton bāad dādī dī sihat sudhrī.} — After the medicine my grandmother's health got better.
  \item \emph{Recall:} say the Punjabi for sweat, then the Punjabi for a dizzy spell, then say \emph{sudharnā} again
\end{itemize}

\begin{tcolorbox}[breakable,colback=teal!4,colframe=teal!35!black,title={Before you move on}]
What does \pa{ਸੁਧਰਨਾ} mean? (\textquotedblleft{}To improve\textquotedblright{}.) Say it once more.
\end{tcolorbox}
\section[\texorpdfstring{\pa{ਵਿਗੜਨਾ} (vigaṛnā)}{ਵਿਗੜਨਾ (vigaṛnā)}]{\pa{ਵਿਗੜਨਾ} (vigaṛnā) — to get worse}

\label{lesson:PA-C325-vigarna}



\begin{quote}
Before the new one: say the Punjabi for a dizzy spell, then the Punjabi for to improve.
\end{quote}

\subsection*{You'll want to know: \pa{ਵਿਗੜਨਾ}}
\textbf{\pa{ਵਿਗੜਨਾ}} — \emph{vigaṛnā} — \textquotedblleft{}to get worse\textquotedblright{}.

\begin{quote}
\pa{ਰਾਤ} \pa{ਨੂੰ} \pa{ਬੱਚੇ} \pa{ਦਾ} \pa{ਬੁਖ਼ਾਰ} \pa{ਵਿਗੜਿਆ।}
\end{quote}

\begin{quote}
\emph{rāt nū̃ bachche dā bukhār vigṛiā.}
\end{quote}

\begin{quote}
At night the child's fever got worse.
\end{quote}

Past (it …): \pa{ਉਹ} \pa{ਮੌਸਮ} \pa{ਸ਼ਾਮ} \pa{ਨੂੰ} \pa{ਵਿਗੜਿਆ}, \emph{oh mausam shām nū̃ vigṛiā}.

\begin{cousinweb}
\textbf{\pa{ਵਿਗੜਨਾ}} (\emph{vigaṛnā}) is inherited from Sanskrit \emph{vighaṭate}, \textquotedblleft{}flies apart, is marred\textquotedblright{}; CDIAL lists Punjabi \emph{vigaṛṇā}, \textquotedblleft{}to be spoilt\textquotedblright{}, under it (CDIAL 11673).
\end{cousinweb}

\begin{grammarlens}[title={What you've built}]
One more word for this chapter.
\end{grammarlens}

\subsection*{Your turn}
\begin{itemize}
\raggedright
  \item \emph{Say it:} \emph{vigaṛnā}
  \item \emph{Say it:} \emph{vigaṛnā}, once more
  \item \emph{Say it:} \emph{rāt nū̃ bachche dā bukhār vigṛiā.} — At night the child's fever got worse.
  \item \emph{Recall:} say the Punjabi for a dizzy spell, then the Punjabi for to improve, then say \emph{vigaṛnā} again
  \item \emph{Write it:} copy \textbf{\pa{ਵਿਗੜਨਾ}} once, checking the whole word
\end{itemize}

\begin{tcolorbox}[breakable,colback=teal!4,colframe=teal!35!black,title={Before you move on}]
What does \pa{ਵਿਗੜਨਾ} mean? (\textquotedblleft{}To get worse\textquotedblright{}.) Say it once more.
\end{tcolorbox}
\section[\texorpdfstring{\pa{ਪਰਹੇਜ਼} (parhez)}{ਪਰਹੇਜ਼ (parhez)}]{\pa{ਪਰਹੇਜ਼} (parhez) — keeping off certain food}

\label{lesson:PA-C325-parhez}



\begin{quote}
Before the new one: say the Punjabi for to improve, then the Punjabi for to get worse.
\end{quote}

\subsection*{You'll want to know: \pa{ਪਰਹੇਜ਼}}
\textbf{\pa{ਪਰਹੇਜ਼}} — \emph{parhez} — \textquotedblleft{}keeping off certain food\textquotedblright{}.

\begin{quote}
\pa{ਡਾਕਟਰ} \pa{ਕਹਿੰਦਾ} \pa{ਹੈ} \pa{ਕਿ} \pa{ਮਿੱਠੇ} \pa{ਤੋਂ} \pa{ਪਰਹੇਜ਼} \pa{ਕਰੋ।}
\end{quote}

\begin{quote}
\emph{ḍākṭar kahindā hai ki miṭṭhe ton parhez karo.}
\end{quote}

\begin{quote}
The doctor says to keep off sweet things.
\end{quote}

A masculine noun.

\begin{cousinweb}
\textbf{\pa{ਪਰਹੇਜ਼}} (\emph{parhez}) is a loan, usually traced to Persian \emph{parhez}, \textquotedblleft{}abstinence\textquotedblright{}.
\end{cousinweb}

\begin{grammarlens}[title={What you've built}]
One more word for this chapter.
\end{grammarlens}

\subsection*{Your turn}
\begin{itemize}
\raggedright
  \item \emph{Say it:} \emph{parhez}
  \item \emph{Say it:} \emph{parhez}, once more
  \item \emph{Say it:} \emph{ḍākṭar kahindā hai ki miṭṭhe ton parhez karo.} — The doctor says to keep off sweet things.
  \item \emph{Recall:} say the Punjabi for to improve, then the Punjabi for to get worse, then say \emph{parhez} again
  \item \emph{Write it:} copy \textbf{\pa{ਪਰਹੇਜ਼}} once, checking the whole word
\end{itemize}

\begin{tcolorbox}[breakable,colback=teal!4,colframe=teal!35!black,title={Before you move on}]
What does \pa{ਪਰਹੇਜ਼} mean? (\textquotedblleft{}Keeping off certain food\textquotedblright{}.) Say it once more.
\end{tcolorbox}
\section[\texorpdfstring{\pa{ਚਬਾਉਣਾ} (chabāuṇā)}{ਚਬਾਉਣਾ (chabāuṇā)}]{\pa{ਚਬਾਉਣਾ} (chabāuṇā) — to chew}

\label{lesson:PA-C325-chabauna}



\begin{quote}
Before the new one: say the Punjabi for to get worse, then the Punjabi for keeping off certain food.
\end{quote}

\subsection*{You'll want to know: \pa{ਚਬਾਉਣਾ}}
\textbf{\pa{ਚਬਾਉਣਾ}} — \emph{chabāuṇā} — \textquotedblleft{}to chew\textquotedblright{}.

\begin{quote}
\pa{ਬੱਚਾ} \pa{ਰੋਟੀ} \pa{ਹੌਲੀ} \pa{ਚਬਾਉਂਦਾ} \pa{ਹੈ।}
\end{quote}

\begin{quote}
\emph{bachchā roṭī haulī chabāundā hai.}
\end{quote}

\begin{quote}
The child chews his bread slowly.
\end{quote}

Past (I …): \pa{ਮੈਂ} \pa{ਰੋਟੀ} \pa{ਹੌਲੀ} \pa{ਚਬਾਈ}, \emph{maiṁ roṭī haulī chabāī}.

\begin{cousinweb}
\textbf{\pa{ਚਬਾਉਣਾ}} (\emph{chabāuṇā}) is inherited from Sanskrit \emph{carvati}, \textquotedblleft{}chews\textquotedblright{}; CDIAL lists Punjabi \emph{cabāuṇā} under it (CDIAL 4711).
\end{cousinweb}

\begin{grammarlens}[title={What you've built}]
One more word for this chapter.
\end{grammarlens}

\subsection*{Your turn}
\begin{itemize}
\raggedright
  \item \emph{Say it:} \emph{chabāuṇā}
  \item \emph{Say it:} \emph{chabāuṇā}, once more
  \item \emph{Say it:} \emph{bachchā roṭī haulī chabāundā hai.} — The child chews his bread slowly.
  \item \emph{Recall:} say the Punjabi for to get worse, then the Punjabi for keeping off certain food, then say \emph{chabāuṇā} again
  \item \emph{Write it:} copy \textbf{\pa{ਚਬਾਉਣਾ}} once, checking the whole word
\end{itemize}

\begin{tcolorbox}[breakable,colback=teal!4,colframe=teal!35!black,title={Before you move on}]
What does \pa{ਚਬਾਉਣਾ} mean? (\textquotedblleft{}To chew\textquotedblright{}.) Say it once more.
\end{tcolorbox}
\section[\texorpdfstring{\pa{ਪੀੜ} (pīṛ)}{ਪੀੜ (pīṛ)}]{\pa{ਪੀੜ} (pīṛ) — an ache}

\label{lesson:PA-C325-pir}



\begin{quote}
Before the new one: say the Punjabi for keeping off certain food, then the Punjabi for to chew.
\end{quote}

\subsection*{You'll want to know: \pa{ਪੀੜ}}
\textbf{\pa{ਪੀੜ}} — \emph{pīṛ} — \textquotedblleft{}an ache\textquotedblright{}.

\begin{quote}
\pa{ਦਵਾਈ} \pa{ਤੋਂ} \pa{ਬਾਅਦ} \pa{ਪੀੜ} \pa{ਘੱਟ} \pa{ਹੈ।}
\end{quote}

\begin{quote}
\emph{davāī ton bāad pīṛ ghaṭṭ hai.}
\end{quote}

\begin{quote}
After the medicine the ache is less.
\end{quote}

A feminine noun.

\begin{cousinweb}
\textbf{\pa{ਪੀੜ}} (\emph{pīṛ}) is inherited from Sanskrit \emph{pīḍā}, \textquotedblleft{}pain, damage\textquotedblright{}; CDIAL gives no Punjabi form, but lists Lahnda \emph{pīṛ}, \textquotedblleft{}pain\textquotedblright{} (CDIAL 8227).
\end{cousinweb}

\begin{grammarlens}[title={What you've built}]
One more word for this chapter.
\end{grammarlens}

\subsection*{Your turn}
\begin{itemize}
\raggedright
  \item \emph{Say it:} \emph{pīṛ}
  \item \emph{Say it:} \emph{pīṛ}, once more
  \item \emph{Say it:} \emph{davāī ton bāad pīṛ ghaṭṭ hai.} — After the medicine the ache is less.
  \item \emph{Recall:} say the Punjabi for keeping off certain food, then the Punjabi for to chew, then say \emph{pīṛ} again
  \item \emph{Write it:} copy \textbf{\pa{ਪੀੜ}} once, checking the whole word
\end{itemize}

\begin{tcolorbox}[breakable,colback=teal!4,colframe=teal!35!black,title={Before you move on}]
What does \pa{ਪੀੜ} mean? (\textquotedblleft{}An ache\textquotedblright{}.) Say it once more.
\end{tcolorbox}
